%% kapa: exam=1 round=32 year=2021 period=1 form=A subject=real_estate subject_ko=부동산학원론 seq=4 number=84
%% source: Q-Net 감정평가사 기출문제 (cst00309) artlSeq=5208397 file=extracted/1008/5208397/A형/2021년도 제32회 감정평가사 1차 1교시 A형.pdf
%% conversion: kordoc 4.24.51 -> tools/gen1.py (auto); review=proofread-sonnet; flags=none
\begin{question}{84}
부동산수요의 가격탄력성에 관한 설명으로 옳지 않은 것은? (단, 다른 조건은 동일함)
\begin{choices}
  \item 수요곡선 기울기의 절댓값이 클수록 수요의 가격탄력성이 작아진다.
  \item 임대주택 수요의 가격탄력성이 1보다 작을 경우 임대료가 상승하면 전체 수입은 증가한다.
  \item 대체재가 많을수록 수요의 가격탄력성이 크다.
  \item 일반적으로 부동산의 용도전환 가능성이 클수록 수요의 가격탄력성이 커진다.
  \item 수요의 가격탄력성이 비탄력적이면 가격의 변화율보다 수요량의 변화율이 더 크다.
\end{choices}
\end{question}
\answer{5}
